\begin{lemma}[One-index near-equal bad-pair counting]
Let $|W_1|=|W_2|=n$ and $|W_3|\in\{n-1,n,n+1\}$. Let $0\leq\varepsilon\leq1$, let $d_{13},d_{23}\leq1$, and suppose $F_{13}[W_1,W_3]$ and $F_{23}[W_2,W_3]$ are $(\varepsilon,d_{13})$- and $(\varepsilon,d_{23})$-super-regular. Call $(x,v)$ bad when its common-neighborhood size is outside $(d_{13}d_{23}\pm2\varepsilon)|W_3|$. If $U\subseteq W_1$ and $W\subseteq W_2$ satisfy $|U|,|W|>2\sqrt\varepsilon n$, then at most a $\sqrt\varepsilon$ proportion of $U\times W$ is bad.
\end{lemma}
\begin{proof}
Put $s=\sqrt\varepsilon$ and $k=|W_3|$. The strict size
assumptions imply $|U|,|W|>0$, since $2sn\geq0$. In particular,
the denominator $|U||W|$ in the asserted proportion is positive.
If $W_3=\varnothing$, every common-neighborhood cardinality is zero.
Both endpoints $(d_{13}d_{23}-2\varepsilon)k$ and
$(d_{13}d_{23}+2\varepsilon)k$ are also zero, so every pair satisfies
the window. There are no bad pairs, and the asserted proportion is
$0\leq s$. Henceforth assume $W_3\ne\varnothing$, so $k>0$.

All three parameters are finite, because they lie between zero and one.
We may therefore do ordinary real arithmetic, interpreting a subtraction
in the nonnegative-reals statements as
$(a-b)_+=\max\{a-b,0\}$. Fix $x\in W_1$, put
$S=N_{F_{13}}(x)\cap W_3$, and write $a=|S|$. The degree clause of
\noderef{SuperRegularPair} gives
\[
 (d_{13}-\varepsilon)_+k\leq a\leq(d_{13}+\varepsilon)k,
 \qquad 0\leq a\leq k.
\]
In particular, the signed real inequalities
$d_{13}k-\varepsilon k\leq a\leq d_{13}k+\varepsilon k$
hold, including when $d_{13}<\varepsilon$.

First suppose $a>\varepsilon k$. For $v\in W_2$, write
$c_v=|N_{F_{23}}(v)\cap S|$, and define the two exceptional sets
\[
 T_- =\{v\in W_2:c_v<(d_{23}-\varepsilon)_+a\},\qquad
 T_+ =\{v\in W_2:c_v>(d_{23}+\varepsilon)a\}.
\]
We show $|T_-|\leq\varepsilon n$. Otherwise
$|T_-|>\varepsilon|W_2|$, and both $T_-$ and $S$ exceed the
thresholds in \noderef{RegularPair}, which applies to $F_{23}$ by
\noderef{SuperRegularPair}. Both sets are nonempty. By the ordered
edge-count definition in \noderef{PairDensity},
\[
 d_{F_{23}}(T_-,S)
 =\frac{\sum_{v\in T_-}c_v}{|T_-|a}
 <(d_{23}-\varepsilon)_+.
\]
Here the strict inequality follows by summing the defining strict
inequalities over the nonempty set $T_-$ and dividing by the positive
number $|T_-|a$. This contradicts the regularity lower bound.
The same reasoning for $T_+$ gives $|T_+|\leq\varepsilon n$:
if its size were larger, its density with $S$ would be strictly greater
than $d_{23}+\varepsilon$, contradicting the regularity upper bound.
These arguments remain valid for $\varepsilon=0$; if a lower-tail
endpoint is zero then that tail is empty because $c_v\geq0$.

For $v\notin T_-\cup T_+$ we have
$(d_{23}-\varepsilon)_+a\leq c_v\leq(d_{23}+\varepsilon)a$.
Together with $a\leq k$, the lower bound implies the following signed
real inequalities:
\[
 \begin{split}
 c_v&\geq d_{23}a-\varepsilon a
      \geq d_{23}a-\varepsilon k\\
    &\geq d_{13}d_{23}k-d_{23}\varepsilon k-\varepsilon k
      \geq(d_{13}d_{23}-2\varepsilon)k.
 \end{split}
\]
Since $c_v\geq0$ as well, this proves
$c_v\geq(d_{13}d_{23}-2\varepsilon)_+k$.
For the upper bound, similarly,
\[
 \begin{split}
 c_v&\leq d_{23}a+\varepsilon a
      \leq d_{23}a+\varepsilon k\\
    &\leq d_{13}d_{23}k+d_{23}\varepsilon k+\varepsilon k
      \leq(d_{13}d_{23}+2\varepsilon)k.
 \end{split}
\]
Both chains use $0\leq d_{23}\leq1$. Also,
$N_{F_{23}}(v)\cap S$ is exactly the common neighborhood in $W_3$.
Thus every bad partner of $x$ belongs to $T_-\cup T_+$, whose size
is at most $2\varepsilon n$.

Now suppose $a\leq\varepsilon k$. Combining this with the degree lower
bound gives $(d_{13}-\varepsilon)_+k\leq\varepsilon k$.
Since $k>0$, cancellation yields
$(d_{13}-\varepsilon)_+\leq\varepsilon$.
If $d_{13}\leq\varepsilon$ then $d_{13}\leq2\varepsilon$ directly;
otherwise the truncated difference is $d_{13}-\varepsilon$, and adding
$\varepsilon$ again gives $d_{13}\leq2\varepsilon$.
Consequently $d_{13}d_{23}\leq d_{13}\leq2\varepsilon$, so the lower
window endpoint is zero. Every $v\in W_2$ has
\[
 0\leq c_v\leq a\leq\varepsilon k
       \leq(d_{13}d_{23}+2\varepsilon)k.
\]
There are no bad partners in this case. We have therefore proved that
every row $x\in W_1$ has at most $2\varepsilon n$ bad partners.

For completeness, the column bound follows by the same two cases, fixing
$v\in W_2$ and taking $S'=N_{F_{23}}(v)\cap W_3$ instead.
Its degree bounds come from \noderef{SuperRegularPair} for $F_{23}$.
If $|S'|>\varepsilon k$, define the lower and upper exceptional sets
in $W_1$ using $(d_{13}-\varepsilon)_+|S'|$ and
$(d_{13}+\varepsilon)|S'|$; \noderef{RegularPair} for $F_{13}$,
with the density sum from \noderef{PairDensity}, bounds each by
$\varepsilon n$. The displayed window chains now have $d_{13}$
in place of the multiplier $d_{23}$ and give the same product window.
If $|S'|\leq\varepsilon k$, cancel the same positive $k$ in its
degree lower bound to obtain $d_{23}\leq2\varepsilon$, and every
intersection lies in the window. Thus every column also has at most
$2\varepsilon n$ bad partners.

Let $B$ denote the number of bad pairs in $U\times W$. Summing the row
bound over $x\in U$ (and restricting each row to $W$) gives
$B\leq2\varepsilon n|U|$. If $s=0$, then $\varepsilon=s^2=0$,
so $B=0$ and the proportion is zero. If $s>0$, then
$\varepsilon=s^2$ and the strict size bound for $W$ gives
\[
 2\varepsilon n= s(2sn)<s|W|.
\]
Dividing $B\leq2\varepsilon n|U|$ by the positive $|U||W|$
therefore gives
$B/(|U||W|)\leq2\varepsilon n/|W|<s$, proving the claimed
non-strict bound. The calculation retains the actual size $k$ throughout;
no equality $k=n$, nor any positive lower bound on $k$ beyond the
explicit nonempty branch, has been used.
\end{proof}
